\chapter{Annexes}
\label{ch:annexes}

\section{Inventaire des réalisations logicielles}

\subsection*{Entrepôt de données hydro-climatiques}

\begin{center}
\begin{tabularx}{\textwidth}{@{}lX@{}}
\toprule
Période & 23 septembre 2025 au 24 août 2026 \\
Volume & $\approx$ \num{7000} lignes Python, $\approx$ \num{1200} lignes SQL de transformation \\
Ingestion & Bibliothèque d'extraction déclarative : pagination, reprise sur erreur,
déduplication par fusion, partitionnement annuel de 1967 à aujourd'hui \\
Transformation & 8 modèles de nettoyage, 3 modèles de rejets tracés, 18 modèles analytiques,
complétés par des tables d'indicateurs calculées en Python \\
Orchestration & 8 traitements planifiés, 3 capteurs événementiels, clés de concurrence par
nature de traitement \\
Stockage & Base relationnelle avec extension de séries temporelles (partitionnement,
compression) et extension géospatiale (jointures spatiales indexées) \\
Infrastructure & 7 services conteneurisés, séparation du code métier et de l'orchestrateur,
volumes de données déclarés hors du cycle de vie des conteneurs \\
\bottomrule
\end{tabularx}
\end{center}

\subsection*{Plateforme d'exploration, de modélisation et de prévision}

\begin{center}
\begin{tabularx}{\textwidth}{@{}lX@{}}
\toprule
Période & 6 novembre 2025 au 24 août 2026 \\
Volume & $\approx$ \num{40000} lignes Python, $\approx$ \num{29000} lignes TypeScript \\
Interface & Application monopage : cartographie, graphiques interactifs, suivi d'entraînement
en temps réel \\
Service & API avec flux d'événements, authentification nominative, journal d'audit, limitation
de débit \\
Cœur métier & Bibliothèque Python sans dépendance d'interface : préparation, modèles métier,
entraînement, explicabilité, contrefactuels, détection d'influence \\
Vérification & \num{552} tests, dont \num{547} passent (mesuré le 24 août 2026) \\
Déploiement & Interface sur infrastructure mutualisée, calcul et bases sur serveur équipé d'un
processeur graphique ; environnements de développement et de production séparés \\
\bottomrule
\end{tabularx}
\end{center}

\section{Sources de données et attributions}

\begin{center}
\begin{tabularx}{\textwidth}{@{}lXl@{}}
\toprule
\textbf{Source} & \textbf{Contenu} & \textbf{Couverture} \\
\midrule
Hub'Eau, piézométrie & Niveaux de nappe aux ouvrages de surveillance & partitions depuis 1967 \\
Hub'Eau, hydrométrie & Débits des cours d'eau (sites, stations, observations) & partitions
depuis 1967 \\
Copernicus ERA5-Land & Réanalyse climatique sur grille \ang{0.1} & depuis 1950 \\
BDLISA & Référentiel hydrogéologique (entités aquifères) & référentiel \\
\bottomrule
\end{tabularx}
\end{center}

\noindent
La profondeur effectivement disponible varie d'une station à l'autre : les partitions
d'ingestion couvrent la période complète, les données retournées par les services web dépendent
de la date de mise en service et de l'exploitation de chaque ouvrage.

\begin{encadre}[title={Mentions obligatoires}]
\emph{Generated using Copernicus Climate Change Service information, 2026. Neither the European
Commission nor ECMWF is responsible for any use that may be made of the Copernicus information
or data it contains.}

\vspace{0.4em}
Source : Hub'Eau, eaufrance.fr \quad\textbullet\quad Source : BDLISA, BRGM et Sandre.

\vspace{0.4em}
Les licences des développements couvrent le code produit, non les données consultées : chaque
source conserve ses propres conditions d'usage.
\end{encadre}

\section{Le cheminement de l'évapotranspiration dans l'entrepôt}
\label{ann:etp}

Une précision s'impose, car la correction ne consiste pas à remplacer une valeur en cours de
route. Les deux grandeurs coexistent, et leurs trajectoires diffèrent.

La variable brute d'ERA5 est \textbf{ingérée telle quelle} en couche Bronze, avec pour seule
opération la conversion des mètres en millimètres. La couche Silver se contente de la typer :
aucun calcul n'y est fait, conformément au principe énoncé au chapitre~\ref{ch:entrepot}.

C'est en couche Gold, dans le modèle de grille mensuelle, que tout se joue. L'évapotranspiration
de référence n'y est pas dérivée de la variable brute : elle est \textbf{calculée à partir d'une
autre source}, les statistiques journalières de température issues du second chemin d'ingestion.
La formule de Hargreaves est appliquée jour par jour, puis cumulée sur le mois. La variable
brute, elle, est conservée dans une colonne distincte de la même table.

\begin{center}
\begin{tabularx}{\textwidth}{@{}llX@{}}
\toprule
\textbf{Couche} & \textbf{Colonne} & \textbf{Contenu} \\
\midrule
Bronze & \code{potential\_evaporation} & Variable ERA5 brute, convertie en millimètres. Aucune
interprétation. \\
Silver & \code{potential\_evaporation} & La même, typée. Toujours aucun calcul. \\
Gold & \code{etp\_totale} & \textbf{Évapotranspiration de référence}, calculée par Hargreaves
depuis les températures journalières. C'est cette colonne qui est consommée. \\
Gold & \code{bilan\_hydrique} & Précipitation moins \code{etp\_totale}. Le SPEI en découle. \\
Gold & \code{etp\_pev\_era5} & La variable brute cumulée, conservée pour la traçabilité et la
comparaison. À ne pas consommer. \\
\bottomrule
\end{tabularx}
\end{center}

Conserver les deux colonnes côte à côte a un intérêt qui dépasse la traçabilité : c'est ce qui a
rendu la mesure possible. Les valeurs de \num{818} et \num{1756} millimètres par an comparées
plus haut sont issues d'une simple requête sur cette table, et non d'un calcul externe. Quiconque
doute de l'arbitrage peut le refaire.

\section{Volumétrie des indicateurs climatiques}

\begin{center}
\begin{tabularx}{\textwidth}{@{}Xr@{}}
\toprule
Mailles de la grille nationale & \num{11496} \\
Valeurs d'indices produites & \num{41960400} \\
Fenêtres d'agrégation & 1, 3, 6 et 12 mois \\
Période de référence & 1991--2020 \\
Échelle de classement & 7 classes \\
Couples maille-mois de la comparaison d'évapotranspiration & \num{30888} \\
Mailles de vérification de l'équivalence des lois & \num{35614} \\
\bottomrule
\end{tabularx}
\end{center}

\section{Glossaire}

\begin{center}
\begin{tabularx}{\textwidth}{@{}lX@{}}
\toprule
\textbf{Terme} & \textbf{Définition} \\
\midrule
Architecture médaillon & Organisation d'un entrepôt en couches de qualité croissante : données
brutes, données nettoyées, données analytiques \\
BDLISA & Référentiel hydrogéologique français décrivant les entités aquifères \\
Contrefactuel & Explication de la forme « quelle modification minimale des entrées changerait
la sortie de telle manière » \\
ETP / ET0 & Évapotranspiration potentielle ; évapotranspiration de référence au sens de la
FAO-56 \\
Hypertable & Table partitionnée dans le temps, permettant compression et interrogation efficace
sur de grands volumes de séries \\
Idempotence & Propriété d'un traitement dont la réexécution ne modifie pas le résultat \\
IPS & Indicateur piézométrique standardisé, position du mois courant par rapport à sa normale
saisonnière \\
L-moments & Statistiques d'ordre linéaires servant à ajuster une loi de probabilité, dont
$\tau_3$ mesure l'asymétrie \\
McKee (classes de) & Échelle à sept classes de sévérité associée aux indices standardisés \\
Modèle TFN & Modèle à fonction de transfert et bruit : le niveau simulé est la somme de
convolutions des forçages par des réponses impulsionnelles calibrées \\
Piézométrie & Mesure du niveau des nappes d'eau souterraine \\
SPI / STI / SPEI & Indices standardisés de précipitation, de température et de bilan hydrique \\
SSFI & Indice standardisé de débit, équivalent de l'IPS pour les cours d'eau \\
STOWA & Standard néerlandais d'évaluation de la qualité des modèles TFN, en quatre critères \\
Valeurs de Shapley & Méthode d'attribution répartissant la contribution des entrées à une
prédiction sur des bases axiomatiques \\
\bottomrule
\end{tabularx}
\end{center}
